\section{Step 5 Complete and store one decoded word}\label{sha-s05}

Load source bytes two and three, place byte two in bits 15 through 8, and combine both with the high fields. Store the completed low word through \texttt{t0}. This is the first observable schedule write. Its bytes in little-endian memory have a different display order from the original block, but loading the numeric schedule word recovers the intended value.

\subsection{Write and inspect}
Replace only the temporary instructions between \texttt{\# BEGIN sch\_low} and \texttt{\# END sch\_low}. Keep the supplied labels outside that region.

\Needspace{10\baselineskip}
\MARSCodeLabel{MIPS assembly}
\begin{minted}[samepage=true]{mars-idea}
    lbu $t3, 2($a0)
    sll $t3, $t3, 8
    or $t2, $t2, $t3
    lbu $t3, 3($a0)
    or $t2, $t2, $t3
    sw $t2, 0($t0)
\end{minted}

Use recipe \texttt{s05-1} with the cumulative program through this edit. The exact recipe names the entry, prepared input and checkpoint. Observe the checkpoint before continuing to its temporary or normal return.

\subsection{Expected observation}
The first schedule word is \texttt{61626380}; later schedule words retain their old contents.

\Needspace{9\baselineskip}
\subsection{Check your understanding}
Why is \texttt{lbu} used for the final source byte?

\begin{itemize}
\item \textbf{A} A signed final byte could set higher fields during OR.
\item \textbf{B} It swaps every pair of bytes.
\item \textbf{C} It writes the schedule automatically.
\end{itemize}
